\scanpage{46}
\mspage{2}
\[
(8)\qquad g_E=\underbrace{g_{V_J(\mathbb R)}}_{\text{tautological}}\oplus\operatorname{sg}(g)\operatorname{id}_\Delta
\]
The action of $\tau$ is given by
\[
(9)\qquad\tau_E=\operatorname{id}_{V_J(\mathbb R)}\oplus(-\operatorname{id}_\Delta)
\]
I set
\[
(10)\qquad\Sigma^*=\Sigma\setminus J=\Sigma-\{R_{\deleted{\unsure{j}}\inserted{i}}\mid i\in I\}
\]
and I am interested in the fundamental groupoid of $\Sigma^*$, as a preliminary exercise in Galois syntax, before looking at \inserted{a little} more complex situations (the equilateral triangle replaced by other configurations rich in symmetries on the sphere, or even on surfaces other than the sphere, or the latter being replaced by the Hopf bundle, with a view to applications to the computation of Teichmüller groupoids...).

It is convenient to introduce a compactification of $\Sigma^*$, so as thereby to have ``base points'' ``close'' to the said $R_i$.
\margin{45}
